% Figure: SIR compartment time series for R0 = 2.5, showing the
% susceptible fall, the infectious peak, and the removed rise, with the
% threshold and the overshoot marked.
\begin{figure}[htbp]
\centering
\begin{tikzpicture}
\begin{axis}[
  width=11cm, height=6.2cm,
  xmin=0, xmax=60, ymin=0, ymax=1.05,
  xlabel={time (arbitrary units)}, ylabel={fraction of population},
  xtick={0,10,20,30,40,50,60},
  ytick={0,0.2,0.4,0.6,0.8,1.0},
  label style={font=\small}, tick label style={font=\scriptsize},
  grid=major, grid style={enggray!25},
  legend style={font=\scriptsize, at={(0.97,0.5)}, anchor=east, draw=enggray!60},
]
% Precomputed SIR curves for R0 = 2.5, gamma = 1/6, S0 ~ 1.
% Susceptible S(t): monotone falling logistic-like curve toward ~0.11.
\addplot[engblue, very thick] coordinates {
  (0,0.999)(4,0.997)(8,0.992)(12,0.980)(16,0.951)(20,0.885)(24,0.755)
  (28,0.560)(32,0.371)(36,0.244)(40,0.176)(44,0.143)(48,0.128)
  (52,0.120)(56,0.116)(60,0.114)};
\addlegendentry{$S/N$ susceptible}
% Infectious I(t): rises to peak ~0.234 near t=33 then falls.
\addplot[engred, very thick] coordinates {
  (0,0.001)(4,0.003)(8,0.008)(12,0.020)(16,0.048)(20,0.109)(24,0.199)
  (28,0.234)(32,0.234)(33,0.234)(36,0.191)(40,0.128)(44,0.076)
  (48,0.042)(52,0.022)(56,0.011)(60,0.005)};
\addlegendentry{$I/N$ infectious}
% Removed R(t): rises to ~0.88.
\addplot[engamber, very thick] coordinates {
  (0,0.000)(4,0.000)(8,0.000)(12,0.000)(16,0.001)(20,0.006)(24,0.046)
  (28,0.206)(32,0.395)(36,0.565)(40,0.696)(44,0.781)(48,0.830)
  (52,0.858)(56,0.873)(60,0.881)};
\addlegendentry{$R/N$ removed}
% threshold line S = 1/R0 = 0.4
\draw[densely dotted, enggray] (axis cs:0,0.4) -- (axis cs:60,0.4);
\node[font=\scriptsize, enggray, anchor=west] at (axis cs:44,0.44) {$S=N/\mathcal{R}_{0}$};
% mark peak time
\draw[densely dashed, enggray] (axis cs:33,0) -- (axis cs:33,0.234);
\node[font=\scriptsize, engred, anchor=south] at (axis cs:33,0.24) {peak};
\end{axis}
\end{tikzpicture}
\caption[SIR compartment time series]{The three compartments of the SIR
model for a basic reproduction number $\mathcal{R}_{0} = 2.5$ in a
wholly susceptible town. The infectious fraction peaks at about
$23\,\%$ exactly when the susceptible fraction falls through the
threshold $1/\mathcal{R}_{0} = 0.40$ (dotted line, dashed peak marker),
after which the infectious count declines because each case now infects
fewer than one successor. The susceptible curve does not fall to zero:
about $11\,\%$ escape, because the infectious pool burns out before
reaching them. The gap between the $60\,\%$ infected at the peak and the
$89\,\%$ infected at the end is the overshoot the outbreak coasts
through on the momentum of the infectious pool.}
\label{fig:vol02:ch07:sir-timeseries}
\end{figure}
